%%=============================================================================
%% Voorwoord
%%=============================================================================

\chapter*{\IfLanguageName{dutch}{Woord vooraf}{Preface}}%
\label{ch:voorwoord}

Ik ben altijd al geïnteresseerd geweest in technologie, en dan niet enkel het schrijven van
code. Vanaf het moment dat ik voor het eerst in aanraking kwam met 3D-modellering,
game development, virtual reality en augmented reality, wist ik dat dit meer was
dan een hobby. Het is een passie, een drive om
dingen te bouwen die mensen kunnen \textit{zien}, \textit{voelen} en
\textit{beleven} in de ruimte om hen heen.

Het was dan ook geen toeval dat ik bij het kiezen van mijn bachelorproefonderwerp bewust
koos voor het domein waar mijn ervaring het kleinst was: augmented reality. Niet
omdat het de gemakkelijkste weg was, maar net omdat ik wist dat de uitdaging me
het meest zou bijleren. Die keuze bleek de juiste te zijn.

Tijdens een gesprek met mijn stageplaats vroeg ik naar inspiratie voor een
bachelorproefonderwerp dat aansloot bij mijn interesses in AR en nieuwe technologie.
Uit dat gesprek groeide een idee dat mij onmiddellijk aansprak: het navigeren op
een begraafplaats via augmented reality. Op het eerste gezicht lijkt dit misschien
een ongewone combinatie, technologie en een plek van stilte en herdenking. Maar
net die tegenstelling maakte het onderwerp voor mij zo boeiend. Hoe kan je een
technologie die doorgaans geassocieerd wordt met entertainment en gaming, inzetten
in een context die vraagt om respect, discretie en menselijkheid?

Dat is de vraag die deze bachelorproef probeert te beantwoorden. Niet enkel
technisch, maar ook menselijk. Want achter elk graf staat een naam, een verhaal,
en iemand die dat graf wil terugvinden.

Deze bachelorproef was een intensief maar ongelooflijk leerrijk traject. Het
ontwikkelen van een native iOS-applicatie met ARKit, het verdiepen in juridische
vraagstukken rond privacy en GDPR, het voeren van een interview met de gemeente
Sint-Gillis-Waas, en het vertalen van al die inzichten naar een werkende proof of
concept: elk onderdeel heeft mij als ontwikkelaar en als onderzoeker doen groeien.

Ik wil graag een aantal mensen bedanken die dit traject mogelijk hebben gemaakt.

In de eerste plaats mijn promotor, \textbf{Dhr. S. Van Impe}, voor de begeleiding
en het vertrouwen doorheen het hele proces. Zijn feedback hielp mij om scherper
te denken en beter te schrijven.

Mijn co-promotor, \textbf{Dhr. L. Gvasalia}, voor de praktijkgerichte inzichten
vanuit de industrie en de inspirerende gesprekken over de mogelijkheden van AR in
professionele contexten.

En tot slot \textbf{Mevr. N. Lesdanon} van de gemeente Sint-Gillis-Waas die de tijd nam om mij
het interne systeem WebBGP te demonstreren en open te praten over de noden en
uitdagingen van begraafplaatsbeheer. Dat gesprek was een van de meest waardevolle
momenten van dit onderzoek.
